\documentclass[a4paper,12pt]{article}
\usepackage[utf8]{inputenc}
\usepackage[T1]{fontenc}
\usepackage[norsk]{babel}
\usepackage{amsmath}
\usepackage{listings}
\usepackage{xcolor}
\usepackage{geometry}
\geometry{margin=2.5cm}
\usepackage{enumitem}
\usepackage[T1]{fontenc}
\usepackage{lmodern}
\usepackage{cmap}
\usepackage{needspace}
\usepackage[colorlinks=true, linkcolor=red!50!black, citecolor=red!50!black, urlcolor=blue]{hyperref}
\usepackage{setspace}
\onehalfspacing

\lstset{
    language=Python,
    basicstyle=\ttfamily\small,
    keywordstyle=\color{blue},
    commentstyle=\color{green!60!black},
    stringstyle=\color{red},
    showstringspaces=false,
    columns=fixed,
    keepspaces=true,
    upquote=true,
    inputencoding=utf8,
    frame=single,
    tabsize=4,
    breaklines=false,
	literate=
     {æ}{{\ae}}1
     {ø}{{\o}}1
     {å}{{\aa}}1
     {Æ}{{\AE}}1
     {Ø}{{\O}}1
     {Å}{{\AA}}1
}

\title{Eulers metode for annenordens differensialligning}
\author{}
\date{}

\begin{document}

\maketitle

\section*{Mål}

Målet med oppgaven er å forstå hvordan vi kan løse en \textit{annenordens differensialligning} numerisk ved å skrive den om til et system av to førsteordens ligninger og bruke Eulers metode.

\section*{Teori: Dempede svingninger}

Vi ser på den klassiske ligningen for en dempet harmonisk oscillator:

\begin{equation}
F = -kx - bv
\end{equation}
Denne kan f.eks. beskrive en pendel som sakte stoppes av luftmotstand. $kx$ er kraften som trekker pendelen mot midten ($x=0$), og $bv$ er luftmotstand. $v$ er selvsagt hastigheten, og vi legger merke til at luftmotstand øker når hastigheten øker.

Siden $F=ma$, $a=\frac{d^2x}{dt^2}$ og $v=\frac{dx}{dt}$ kan kraftligningen skrives som differensialligningen
\begin{equation}
\label{eq:diff}
\frac{d^2x}{dt^2} =- \alpha \frac{dx}{dt} - \beta x,
\end{equation}
der $\alpha=b/m$ og $\beta=k/m$.

Dette er en annenordens differensialligningligning siden den involverer den andrederiverte. For å bruke Eulers metode må vi skrive den om til to førsteordens ligninger.

Vi gjeninnfører

\[
v = \frac{dx}{dt},
\]

og omskriver differensialligningen \eqref{eq:diff} som:

\begin{subequations}
\begin{align}
\frac{dx}{dt} = v &\iff dx = v \, dt \label{eq:dx}\\
\frac{dv}{dt} = -\alpha v - \beta x &\iff dv = -(\alpha v +\beta x) \, dt \label{eq:dv}
\end{align}
\end{subequations}

Nå kan vi bruke Eulers metode på begge ligningene samtidig:

\begin{subequations}
\begin{align}
t_{n+1} &= t_n + dt,\\
x_{n+1} &= x_n + dx \\
\implies x_{n+1} &= x_n + v_n\, dt &(\text{fra \eqref{eq:dx}}), \\
v_{n+1} &= v_n + dv\\
\implies v_{n+1} &= v_n - (\alpha v_n + \beta x_n)\, dt &(\text{fra \eqref{eq:dv}}).
\end{align}
\end{subequations}

\section*{Oppgave 1: Forstå systemet}

La parameterne være gitt ved

\[
\alpha =  0.1, \quad \beta = 1
\]
og la initialbetingelsene være

\[
x(0) = 1, \quad v(0) = 0.
\]
Med andre ord har vi $(t_0,v_0,x_0) = (0,0,1)$.


Bruk $dt = 0.1$ og regn ut følgende.

\begin{enumerate}[label=\alph*)]
\item Regn ut $t_1$, $v_1$ og $x_1$.\footnote{$(t_1,v_1,x_1) = (0.1, -0.1, 1)$.}
\item Regn ut $x_2$ og $v_2$.\footnote{$(t_2,v_2,x_2) = (0.2, -0.195, 0.99)$.}
\item Fortsett til du går lei.
\end{enumerate}

\newpage

\section*{Oppgave 2: Programmer funksjonen}
Programmer funksjonen og definer parameterne:
\begin{lstlisting}
x0 = 1
v0 = 0
alpha = 0.5
beta = 4

t0 = 0
t_slutt = 10
dt = 0.1

def f(x, v):
    return -alpha*v - beta*x
\end{lstlisting}

\section*{Oppgave 3: Implementer Eulers metode}

\begin{lstlisting}
x = x0
v = v0
t = t0

while t <= t_slutt:
    x = x + v * dt
    v = v + f(x, v) * dt
    t = t + dt
    print(t, v, x)
\end{lstlisting}

\begin{enumerate}[label=\alph*)]
\item Hvilken linje tilsvarer formelen for $v_{n+1}$?
\item Sjekk at utregningene fra oppgave 1 stemmer.
\end{enumerate}

\section*{Oppgave 5: Tegn grafen}

\begin{lstlisting}
import matplotlib.pyplot as plt

x_values = [x0]
v_values = [v0]
t_values = [t0]

x = x0
v = v0
t = t0

while t <= t_slutt:
	t = t + dt
    x = x + v * dt
    v = v + f(x, v) * dt
    
    
    x_values.append(x)
    v_values.append(v)
    t_values.append(t)

plt.plot(t_values, x_values)
plt.xlabel("Tid")
plt.ylabel("Posisjon")
plt.title("Dempede svingninger med Eulers metode")
plt.grid()
plt.show()
\end{lstlisting}

\section*{Oppgave 6: Sammenlign med eksakt løsning}
Differensialligningen har tre forskjellige løsninger. Den første for $\alpha^2 < 4\beta$:

\begin{equation}
x(t)=
\left[
x_0\cos\!\left(\frac{\sqrt{4\beta-\alpha^2}}{2}\,t\right) +
\frac{2v_0+\alpha x_0}{\sqrt{4\beta-\alpha^2}}
\sin\!\left(\frac{\sqrt{4\beta-\alpha^2}}{2}\,t\right)
\right] e^{-\frac{\alpha}{2}t}
\end{equation}
Den andre for $\alpha^2 = 4 \beta$:

\begin{equation}
x(t)=
\left[
x_0+\left(v_0+\frac{\alpha}{2}x_0\right)t
\right]
e^{-\frac{\alpha}{2}t}
\end{equation}
Og den tredje for $\alpha^2 > 4 \beta$:

\begin{equation}
x(t)=
\frac{v_0-r_- x_0}{r_+ - r_-}\,e^{r_+ t}
+
\frac{r_+ x_0-v_0}{r_+ - r_-}\,e^{r_- t},
\end{equation}
der
\[
r_{\pm}=\frac{-\alpha \pm \sqrt{\alpha^2-4\beta}}{2}.
\]

I Python ser dette slik ut:
\begin{lstlisting}
import numpy as np

def X(t):
    D = alpha**2 - 4*beta

    # Underdempet (komplekse røtter)
    if D < 0:
        omega = np.sqrt(4*beta - alpha**2) / 2
        C1 = x0
        C2 = (v0 + (alpha/2)*x0) / omega
        return np.exp(-alpha*t/2) * (
            C1*np.cos(omega*t) + C2*np.sin(omega*t)
        )

    # Kritisk dempet (dobbel rot)
    elif D == 0:
        r = -alpha/2
        C1 = x0
        C2 = v0 - r*x0
        return (C1 + C2*t) * np.exp(r*t)

    # Overdempet (to reelle røtter)
    else:
        r1 = (-alpha + np.sqrt(D)) / 2
        r2 = (-alpha - np.sqrt(D)) / 2
        
        C2 = (v0 - r1*x0) / (r2 - r1)
        C1 = x0 - C2
        
        return C1*np.exp(r1*t) + C2*np.exp(r2*t)

T = np.linspace(t0, t_slutt, 500)

plt.plot(T, X(T), "--")
\end{lstlisting}


\end{document}